% Author: Maxwell Chen
\pgfplotsset{width=7cm,compat=1.16}

\qns{Gram-Schmidt Process}

Given a vector space $V$ we have seen that there are different ways to represent vectors in this space through the use of a \textbf{basis.} Consider an arbitrary set $S = \{ \vec{v}_{1}, \cdots \vec{v}_{k} \}$ that is a basis for $V.$ It would be easier to work with a set of vectors that are mutually orthogonal and of unit length. Recall that vectors are \textbf{orthogonal} when their inner product is zero. Geometrically the vectors will form a $90^{\circ}$ angle.

The term \textbf{orthonormal} collapses these two terms into one: a set of vectors is orthonormal \textbf{if and only if}:
  \begin{itemize}
    \item Any pair of vectors is orthogonal (for vectors $\vec{v}_{i}, \vec{v}_{j} \in S$ where $\vec{v}_{i} \neq \vec{v}_{j}$, the inner product $\innp{\vec{v}_{i}}{\vec{v}_{j}} = 0.$)
    \item Each vector is of unit length (for all vectors $\vec{v}_{i} \in $ $S$, $\norm{\vec{v}_{i}} = 1$).
  \end{itemize}

The Gram-Schmidt process is a method that allows us to take an arbitrary basis $S$ for some vector space $V$ and create a basis $U = \{ \vec{u}_{1}, \cdots \vec{u}_{n} \}$ that spans the exact same space, yet is now orthonormal.

\begin{enumerate}

\qitem Show that span($\vec{v_1}, \vec{v_2}$) is the same as span($\vec{v_1}, \vec{v_2} - \alpha\vec{v_1}$), where $\alpha$ is some arbitrary constant.

\sol {
  The definition of span is all vectors that can be "reached" using a linear combination of a set of vectors. Therefore,
  $$\text{span}(\vec{v_1}, \vec{v_2}) = \{\vec{v} : \vec{v} = \beta_{1} \vec{v}_{1} + \beta_{2} \vec{v}_{2} \} \ \ \text{for arbitrary constants} \ \beta_{1}, \beta_{2}.$$ 
  Following a similar train of thought,
  $$\begin{aligned} \operatorname{span}\left(\vec{v}_{1}, \vec{v}_{2}-\alpha \vec{v}_{1}\right) 
  &= \{\vec{v} : \vec{v} = \beta_{3} \vec{v}_{1}+\beta_{4}\left(\vec{v}_{2}-\alpha \vec{v}_{1}\right) \} \\ 
  &= \{\vec{v} : \vec{v} = \beta_{3} \vec{v}_{1}+\beta_{4} \vec{v}_{2}-\beta_{4} \alpha \vec{v}_{1} \} \\ 
  &= \{\vec{v} : \vec{v} = \left(\beta_{3}-\beta_{4} \alpha\right) \vec{v}_{1}+\beta_{4} \vec{v}_{2} \} \ \ \text{for arbitrary constants} \ \beta_{3}, \beta_{4}.
  \end{aligned}$$
  All $\alpha$ and $\beta$ are arbitrary constants that we can set. If we set $\beta_4 = \beta_2$ and $\beta_3 = \beta_1 + \alpha\beta_2$, then these spans are the same!
}

\qitem Let us iteratively generate the orthonormal set $U$ from $S$. Start by arbitrarily picking a vector $\vec{v}_{1}$ from $S.$ How can we create a new vector $\vec{u}_{1}$ for $U$ so that it spans the same space as the original vector, yet ensures that orthonormality is preserved for our new set?

\sol {
  Since we are only taking into consideration a single vector, we do not need to worry about orthogonality. We do however, need to make sure its norm is $1.$ So,
  $$\vec{u}_{1} = \frac{\vec{v}_{1}}{\norm{\vec{v}_{1}}}$$
}

\qitem Now, consider an arbitrary vector $\vec{v}_{2}$ in $S$. What is the projection of $\vec{v}_{2}$ onto $\vec{u}_{1}?$

\sol {
  The projection of an arbitrary vector $\vec{y}$ onto another vector $\vec{x}$ is:
  \begin{equation}
    \text{proj}_{\vec{x}} \vec{y} = \frac{\innp{\vec{y}}{\vec{x}}}{\innp{\vec{x}}{\vec{x}}} \vec{x}
  \end{equation}
  This can be seen through the following visual:
  \begin{center}
    \includegraphics[width=2in]{\bank/inner-products/figures/projection_visual.png}
  \end{center}
  The cosine formula for inner products is:
  \begin{equation}
    \cos(\theta) = \frac{\innp{\vec{x}}{\vec{y}}}{\norm{\vec{x}} \norm{\vec{y}}}
  \end{equation}
  Substitute for $\cos(\theta),$ and using a unit vector for $\vec{x}, \frac{\vec{x}}{\norm{\vec{x}}}$ we get the projection formula above. 

  Therefore, if we project $\vec{v}_{2}$ onto $\vec{u}_{1},$ since $\vec{u}_{1}$ has unit norm, the projection will be:
  \begin{equation}
    \text{proj}_{\vec{u}_{1}} \vec{v}_{2} = \innp{\vec{v}_{2}}{\vec{u}_{1}} \vec{u}_{1}
  \end{equation}
}

\qitem How can we create $\vec{u}_{2}$ for our orthonormal set using the projection defined in the previous part? \\
(Hint: Take a look at part (a))

\sol {
  Now, we must consider orthogonality as well since we have two vectors in our set. For two vectors $\vec{u}_{1}, \vec{u}_{2}$ to be orthogonal, their inner product must to be 0. In other words, this means that there is nothing--no trace, no component--of $\vec{u}_{1} \text { in } \vec{u}_{2}$, and vice versa. 
  We also know from part (a) that we can make $\vec{u}_{2}$ a linear combination of $\vec{u}_{1}$ and $\vec{v}_{2}$ without changing the span. In order to make $\vec{u}_{2}$ orthogonal to $\vec{u}_{1},$ we will subtract the projection of $\vec{v}_{2}$ onto $\vec{u}_{1}$ by doing the following:
  $$\vec{q_2} = \vec{v_2} - \innp{\vec{v}_{2}}{\vec{u_1}} \vec{u_1}$$
  Geometrically, it will look like the following:
  \begin{center}
    \includegraphics[width=2in]{\bank/inner-products/figures/gram_schmidt_visual.png}
  \end{center}

  By subtracting the projection of $\vec{v}_{2}$ onto $\vec{u}_{1}$, we are removing the $\vec{u}_{1}$ component of $\vec{v}_{2}$ from $\vec{v}_{2}$. So, $\vec{q}_{2}$ is now orthogonal to $\vec{u}_{1}$, so mutual orthogonality is preserved for all vectors in our new set $U$. 
  Lastly, we make sure that $\vec{u}_{2}$ is normalized, so:
  $$\vec{u}_{2} = \frac{\vec{q}_{2}}{\norm{\vec{q}_{2}}}$$
}

\qitem How can we use the information above to develop the process of Gram-Schmidt for all future vectors $\vec{v}_{3}, \cdots, \vec{v}_{n}?$ Write out the process for creating an arbitrary vector $\vec{u}_{i}$ for our orthonormal set $U$.

\sol {
  If we try repeating the process for $\vec{v}_{3},$ we would have to make $\vec{u}_{3}$ orthogonal to both $\vec{u}_{1}$ and $\vec{u}_{2}.$
  We can again do this by taking the projection of $\vec{3}$ onto both $\vec{u}_{1}$ and $\vec{u}_{2}$ simultaneously, and subtract the individual projections to make $\vec{q}_{3}$ orthogonal to both $\vec{u}_{1}$ and $\vec{u}_{2}.$
  Define the following vector:
  $$\vec{q}_{3} = \vec{v}_{3} - \innp{\vec{v}_{3}}{\vec{u}_{1}} \vec{u}_{1} - \innp{\vec{v}_{3}}{\vec{u}_{2}} \vec{u}_{2}$$
  Then we can normalize accordingly to make sure $\vec{u}_{3}$ has norm $1.$
  $$\vec{u}_{3} = \frac{\vec{q}_{3}}{\norm{\vec{q}_{3}}}$$
  In general, we can do this for every vector by subtracting the component of every previous vector already in $U$, ensuring that our new vector $\vec{q}_{i}$ is orthogonal to the rest. 
  $$\vec{q}_{i} = \vec{v}_{i} - \sum\limits_{j = 1}^{i - 1} \innp{\vec{v}_{j}}{\vec{u}_{j}}$$
  And remember to normalize $\vec{q}_{i}.$
   $$\vec{u}_{i} = \frac{\vec{q}_{i}}{\norm{\vec{q}_{i}}}$$
  We can repeat this for all vectors in $S$ to generate a new, orthonormal set $U$ that we can use in place of $S$ with orthonormal vectors completing our algorithm.
}
\end{enumerate}